\chapter{Theoretical reformulated frameworks}

Since the previous analysis of the learning model, we can then notice a few disadvantages and a few very subtle points that we omitted. It is apparent that the analysis is somewhat entangled, with various situations, interpretations, setting and assumptions together, and with a non-unified structure of models in question. The system in which the model is considered and analysed is also missing in essence. There are a lot of ambiguities around what is considered and what is not, even with probabilistic nature and the deterministic aspect. Terms like \textbf{latent spaces} and many do not capture the entire dynamic or meaning of the system. It is then we introduce our novel analytical framework, namely the \textbf{Unified Unitary Construct} framework. It is however, can be said otherwise that the learning theory is rather not so novel, but approaching the problem is a more distinct way of the already existing framework. 
\section{Motivation}
The predecessors of neural network, for operations and purposes, take in the form of a specific application of function constructs, finite singleton algorithms. For such purposes, a specific classical model $m$ takes the form of a single structure, with complex construction, often offering fixed architecture of its representation. Neural network is an attempt, for its philosophy, to change the aforementioned philosophy of constructing machine model by considering compute unit \cite{mcculloch_logical_1943,Rosenblatt1958ThePA}. 

Motivated by earlier analysis and observations, it is apparent that the analysis is entangled, with various situations, interpretations, setting and assumptions together, and with a non-unified structure of models in question. The system in which the model is considered and analyzed is also ineffective when considering larger structures or more `realistic' consideration aside from very simplified, strong system itself. There are a lot of ambiguities around what is considered and what is not, even with probabilistic nature and the deterministic aspect. Terms like \textbf{latent spaces} and many do not capture the entire dynamic or meaning of the system. It is then we introduce our novel analytical framework, namely the \textbf{Unified Unitary Construct} framework. It is however, can be said otherwise that the learning theory is rather not so novel, but approaching the problem is a more distinct way of the already existing framework. 

Recalled the classical learning framework, for $h$. We know and construct the principle of the system by settle on $h:\mathcal{X}\to \mathcal{Y}$, for any arbitrary defined in-out 2-tuple. Hence, the overall construction of a given model can be abstracted into layers of responsive results. This in particular draws some advantages, for example, can be easily approximated or expressed in the black box model style, or empiricism approach. However, it also comes with the drawback that we have seen, in the sense that the interpretation of such black box model is very narrow, even after we add different kind of flavours, for example, the assumptions of the Gaussian noise. This is majorly because of applying and addition of different components without a general framework, leading to a lot of messiness when looking directly into the theory itself. Specifically, for example, is the fact that there exists two kinds of learning analysis, one considers the information-theoretic approach of information exchange, while the other is simply the classical approach given in such the intrinsic noise error given by the observation set itself, keeping the classical theory intact. How to formalise this is a big hurdle to incorporate each other in, without considering the wider framework - one of the main problem with the speciality of learning theory. Nevertheless, at least accordance to general theory of the current date, formalization in to unit-like processing unit as \textit{neural network components} is one thing that should be considered for the learning framework. This is also one of the main point for the word \textbf{Unified} in our theory. 


%\subsection{Probability}
%
%Probability is a loose concept, most generally speaking. Usually, it refers to the uncertainty in our action, hence %taken of virtue of subjectivity given certainty being not absolute. As reconciled by \cite{Nau2001_DeFinettiWasRight,%deFinetti1990_TheoryProbability,Elliot2019_DeFinettiPhilosophy}, probability can be described to be of the \textit%{degree of belief} of certain agent abstraction - an object or specimen that receives and act on the environment it %beholds. In such sense, we can also interpret our system to explain the uncertainty that unfolds. 
%
%Taken as an observer tasked with predicting and interpreting the system $\mathcal{S}_{0}$, then, by the standard %definition of ambiance space. A system that is \textit{total predictable} and \textit{deterministic} if, the encoding %space is empty, $\mathcal{E}=\varnothing$, and $\mathcal{L}^{*}=\mathcal{L}_{E}$. What this means is that all %descriptions are known, the language is observable from such, and there are no encoding. This is the \textit{ideal %system}. 

\section{The concept of learning}

Within the authoritative sources on the modern deep learning theory, like \cite{neyshabur2015norm,neyshabur2017exploring,NJacot:2018:NTK,arora2019finegrained,lee2019wide,chizat2019lazy,zou2019implicit,bartlett2021deep,roberts2022principles,wei2022theory,yang2019scaling} in comparison, the new theory of modern networks, or deep learning, remains a very troublesome regime with a lot of unspecified behaviours or outcomes leading on its structures. Unsolved problems emerged from the structural usage, for example, double descent, grokking, neural scaling limits, etc, all of which hinders and expose the immaturity of the concurrent theory. As such, a direction to take on to resolve such issue must be contained, effectively nonetheless. 

\section{Neurons perspective}

The formalization to be \textit{atoms} comes from the biological inspiration of exactly the biological neurons in human and species' brain. However, the mechanism and how they operate is a bit more interesting. 

For the biological more advanced equivalent, the brain consists of a very large number of highly connected elements, called \textit{neurons}. The simplified structure of such neuron consists of three components: the \textit{dendrite}, which is tree-like receptive network that carry signals (electrical) to the cell body; the \textit{cell body} effectively processes these signals, usually as sums and threshold response; and the \textit{axon}, a single long fibre that carries signal from the cell body to other neurons. The point of contact between an axon and a dendrite of another cell is called a synapse. The neural network, hence, is established by the synapses and various arrangements of the neuron. 

This configuration of biological neuron is particularly powerful. Not only that they can process information allowing someone to write this text, but also the fact that of billions or neuron, they parallelly work at the same time. Furthermore, the intricacy lies further ahead than the above formation - the entire structure is of too much complexity, for the brain itself. However, a principle can be seen - they are contained of \textit{building blocks} together, in which case is neuron. And \textbf{artificial networks} also take this approach. 

\begin{note}
From such, we can conclude that, we have to have some classification that requires the smallest component neuron to be the \textit{atoms} that everything else is formed upon.
\end{note}

One of the most fundamental thing of the structure of the \textit{neural network formalism}, hence, is the notion of a \textit{unit neuron}. Specifically, a neuron $x$ is defined, and designed to be a discrete operating unit on its own. The word \textit{operating} here means that it has all the facility requires for an input-output process - the most simple one includes the input receiver, the processing formulae, and the output transmitter. 

The neural network formalism is best expressed by the following principle. This theorem is by design, and thus our design is to fix the foundational minimization to a specific neuronal basis module. Such is to say, we state the principle for \textit{inductive recursion} of neuronal system collapse. 

\begin{theorem}[Fundamental theorem of neural formalism]\label{thm:fundamental_neuron}
Let $\mathcal{S}$ be an automated operational system with internal composition system, admitting
\begin{enumerate}[topsep=1pt,itemsep=1pt]
    \item nontrivial input-output interaction,
    \item internal closure under composition, and
    \item irreducibility with respect to its operational primitives.
\end{enumerate}
Then any minimal operational component of $\mathcal{S}$ must admit a representation whose cardinality and arbitrary (analytic, topological, etc.) structures are equivalent to that of the minimally defined neural structure $\mathfrak{N}(\mathcal{N},\mathcal{C})$.
\end{theorem}

We begin by examining the model of a typical neuron. Note that, this is the simplest one, but also is the fundamental block. We have the following. 

\begin{definition}[Neuron]
    Given the neural network formalism. Define a construct $x$. Then $x$ is called a \textbf{neuron} if it belongs component-wise to the class $\mathcal{N}$ of all neurons, which is minimally expressed by $\mathcal{N}(n_{i},n_{o},M(\dots))$, where $n_{i}$ is the input handler, $n_{o}$ is the output handler, and $M(\dots)$ is the internal system. %\index{classical neuron}
\end{definition}

From this, we then come up with the definition of the minimally defined neuron - the single most iteration of the above definition. 

\begin{definition}[Minimally defined neuron]
    Given the neural network formalism. Then, for $x\in \mathcal{N}$, we call a neuron \textbf{minimally defined} if it belongs to the class $\mathcal{N}_{0}\subset \mathcal{N}$ of $\mathcal{N}_{0}(n_{i}[\mathbf{w}],n_{o}[\mathbf{w}'],M(f,b))$ for $f:n_{i}[\mathbf{w}]\times b \to \mathcal{O}$ is the internal function of $M$. $\mathcal{O}$ is the domain of $n_{o}$ in which it receives the value, and in cases, there exists no $\mathbf{w}'$ configured for the output channel. %\index{minimally defined classical neuron}
\end{definition}

The notation $n_{i}[\mathbf{w}], n_{o}[\mathbf{w}']$ inherently indicate the notion of \textit{control} of the neuron on the two gate of input and output. More specifically, a given neuron, aside from the upper-level expression of $\mathcal{N}(\dots)$, can also be realized by its \textit{parameters}, or rather, its configuration of the neuron itself. By that, for $x \in \mathcal{N}$, the \textit{configuration parameters} for the minimally defined neuron can be taken in the form $\mathcal{C}(\mathbf{w}_{i}, \mathbf{w}_{o}, \mathbf{w}_{M})$, here we use instead the pairing notation. The letter $\mathbf{w}$ used here is historical by certain accounts: the original idea comes from the term \textbf{weight}, in which the neuron can conceptually influence the input, for example, the signal received or information received, by certain amount, either \textit{downplay} it or \textit{signify} it.

We clearly clarify the need for both $\mathcal{N}$ and $\mathcal{C}$ here. We know, that we want the neuron to be an \textit{operating unit}. This means that for it to be fully realized, it needs to also operate, and hence, to be 'observed'. Hence, you need both the description of the parameters which defines it - taken the interpretation where each parameter and configuration is one specific building block on its own, then the block of all those parameters form the shape of the neuron. Similarly, the \textit{operations} on such neuron assure the interaction and working mechanism of all those components together, inside the neuron, by itself. Hence, to fully, minimally express a minimalized neuron, you need to have both $\mathcal{N}$  and $\mathcal{C}$ as its minimality requirement. We then revise it to the following definition. 

\begin{definition}[Classical neuron class] \label{def:neuclass}
    Given the \textit{neural network formalism}. Then, we define a \textbf{neuron} $A$ to be \textit{minimally defined} if it is equivalent to any unit $x$ of the class $\mathfrak{N}_{m}(\mathcal{N}, \mathcal{C})$, called a \textit{neuron class}, where the 2-tuple expanded to $\mathcal{N}(n_{i}, n_{o}, M)$ and $\mathcal{C}(\mathbf{w}_{i}, \mathbf{w}_{o}, \mathbf{w}_{M})$. %\index{classical neuron class}
\end{definition}

\begin{figure}[!ht]
    \centering
    \resizebox{0.6\textwidth}{!}{%
    \begin{circuitikz}
    \tikzstyle{every node}=[font=\Large]
    
    % Manually place the label above the box
    \node at (29,7) {\Large $x\in\mathfrak{N}_{m}(\mathcal{N}, \mathcal{C})$};
    
    % Outer dashed box
    \draw[dashed] (21.75,6.25) rectangle (36.25,-1.25);
    
    % Inner boxes
    \draw (22,6) rectangle node {\Large $n_{i}[\mathbf{w}]$} (23.75,-1);
    \draw (36,6) rectangle node {\Large $n_{o}[\mathbf{w}']$} (34.25,-1);
    \draw (24,6) rectangle node {\Large $M(\dots)$} (34,-1);
    \end{circuitikz}
    }%
    \caption{Minimal neuron structure}
    \label{fig:lable}
\end{figure}

With this, we formalize the abstract neuron into two descriptions. Note that, however, the \textit{minimally defined} neuron has a parameterized exposure protocol to the construction. In practice, for any neuron $a \in \mathfrak{N}$, the tuple $(\mathcal{N}, \mathcal{C})$ can be different, not to say complex. One of the strong points of this construction of the neural class is that it is \textit{recursive}. The following proposition might help aid in understanding why it is recursive.

\begin{proposition}[Compound construction]
    An abstract neuron $A$  can be dissected into three most important compartments $(\mathrm{I},\mathrm{O},M)$, where $\mathrm{I},\mathrm{O}$ is respectively the in-out interface. Hence, a cluster of neurons $\{A_{i}\}_{i \leq n}$ for such $\lvert \{ A_{i} \} \rvert \geq 3$ can be compartmentalized into a new neuron unit, that is, $\{ A_{i} \}\in \mathfrak{N}$, if and only if $\lvert \{ A_{i} \} \rvert \geq 3$. 
\end{proposition}

\begin{proof}
    This mostly comes of as definition $(2.1)$. Notice that for any structure of a neuron to be minimally defined, the tuple $\mathcal{N}$ must minimally contain $(n_{i},n_{o},M)$. We want them to be discretely defined, then a cluster of neuron will have to have at least one neuron such that to satisfy the requirement of the 3-tuple component. If $n=\lvert \{ A_{i} \} \rvert=1$, it reduces to a singular neuron; for $n=2$, at least one compartment do not have any component, hence minimally we need 3 neurons to effectively be arranged as a neuron unit.
\end{proof}

Hence, we can recursively define neuron, either going up or down, if they satisfy such condition. However, usually, we will only go up, as we will have to define in detail the set $\{ A_{M,i} \}$ of all possible \textit{minimally defined} neuron. \footnote{At this point, I am considering using \textit{group theory} to aid such set compartment.}

There are several assumptions and properties accompanied by this type of neuron structure that we would like to note.
\begin{enumerate}[itemsep=1pt,topsep=1pt]
    \item This neuron structure works \textit{sequentially}. A minimally defined neuron, such as in a practical setting, will have $\mathcal{C}$ to support sequential operation. What this means is present in the definition of the minimal neuron, such that $\mathcal{C} = (\mathbf{w}_{i}, \mathbf{w}_{o}, \mathbf{w}_{M})$, and the operation follows $i \to M \to o$ for an input-process-output session.
    
    \item By extension of the point above, the neuron order in a cluster will also be \textit{fixed}. Usually, this follows a single dimension, or, for example, if there exists a singular span directional vector $\vec{dr}$ that indicates the direction of operation in the cluster.
    
    \item We suppose the \textit{independence} of any components inside a neuron. This is meant for both \textit{operational independence} for non-connective and resultant actions case. If subcomponents are connected in a causal form and so on, then such independence is violated, but not otherwise. 
    
    \item The neuron is expressed in a \textit{parameterized manner}. That is, the behaviour of the neuron can be expressed, controlled, and designed by modifying and constructing a set of non-realized parameters. If this set of parameters is finite, we say that it is \textit{closed parameterization}. If not—i.e., there can be infinitely many parameters—we call it \textit{open parameterization}.
    
    \item This neuron structure is called the \textit{forgetful neuron}, simply because it works, in principle and in actuality, as a processing unit, in the minimally defined case. We will soon touch upon this definition for the detailed structure of the minimally defined neuron.
\end{enumerate}

When we structure our neuron in this way, it is natural to ask whether the neural unit can be expressed similarly to a \textit{finite automaton} or not. But aside from such an analogy, our minimal neuron is already very powerful. Certain clusters of such neurons, embedded with structure $\mathcal{L}$ of \textit{layer-order clustering} (which can be thought of as a more formal notion for a layered neural network), can approximate any continuous representation scheme $c: \mathcal{X} \to [\cdot, \Sigma]$, given an arbitrary closed interval and to an arbitrary degree of accuracy. We often refer to this as the \textit{universal approximation theorem} for the minimally defined neuron system, and it will serve as one of the focuses of the future analysis.

\subsection{Analysis}
The above construction of the fundamental neuron fundamental create the component model of the minimum processing unit which will be used in various collections of system of related mean. In mathematical modelling context, the neuron model is \textit{descriptive}, meaning that it aims specifically not only for the input-output protocol, but also the unit neuron itself. To do this, we recall any given concrete, functional structure under mathematical formalism is given of the basis, on which abstraction and quantification (without physical realization). Hence, our model here aligns well with such mathematical formalism, in the sense that they end up with abstraction without the impromptu need to signify the inner physical realization - or existence itself. And doing this leaves us with formalizing the \textit{minimally defined neuron} of the class $\mathfrak{N}(\mathcal{N},\mathcal{C})$ into specific requirement. 

The inspiration for the concept of a neuron is a biological one: it stems from the study of the brain or any given processing organs of species with specialized region to be called brain. Considering such, we found out that the brain mostly consists of the fundamental component called biological neuron and its supporting structures - other resource matters of the inner brain that helps to make it function. Without physical realization depends on many things, but most of the time, it is expressed by the proposition that is to remove the physical attachment from any given object of interest, giving it a more quantified look. For a neuron, it can be treated in such way by formalize the neuron component without its supporting structure, replacing electrical signal and other innate signalling mechanism with abstract 'data flow' and 'operations', without much worry about the actual physical realization of chemical reaction and channel - as we do not need such. 

Back there, we all said of the neuron class. What does this mean? For starter, this section on the neuron is the relatively minimal construction we can get it onto. Recall that we say mathematical formalism is to put objects in quantifiable notions, and abstractions of the given subject; also, to express them in a language (or framework, depends on your choice of lexical sense) that support such. Hence, one of the first thing we would like to do, is to formalize the neuron into specific \textit{description scheme}; that is, we would like to treat our objects - the neurons - specifically in certain descriptions that captures the essence of the neural unit, without much ambiguity, and an overhead view of the neuron. The tuple $(\mathcal{N},\mathcal{C})$ did exactly that - it tells us that the neuron is expressed by that separation of descriptions - in the language of mathematical quantification; for then, $\mathcal{N}$ stands for the quantitative \textit{parameters} that characterize the resources, component scale, or the \textbf{mass} of a given neuron, which $\mathcal{C}$ stands for the operations that governs the neuron's internal dynamics. One feasible assumption here is that the neuron does not exhibit any of its operation, that is, interfering with the macroscopic world outside itself, aside from the input-output formalism. Those two descriptions, one governs the mass, one govern the internal mechanics that confounds the subject's behaviours, is what the mathematical description do - and it works very well as it is. 
\vspace{2mm}

In such formalism, $\mathcal{N}$ is typically easier to specify than $\mathcal{C}$ - which is trivial since the internal mechanics is harder to define and construct, for there exists many relationship and connection between components in the same system. And, considering that for any working system, especially for an interactive one, which inherently dependent on the scale $t$, or given any reference point of operation to be referred upon, there can be either infinite configuration, or infinite type of ordering of the system itself. We may come to the analysis on \textit{invariant structure} of the space of configuration, but otherwise, it is uncommon to have invariant structure presented in the majority inside a typical system, except fitting of certain criteria.  

Overall, the neuron formalism of capturing them into a neural class also helps to figure out the \textit{macroscopic configuration} and the \textit{microscopic configuration} of any given neuron. Under a certain instance of the neuron class $\mathfrak{N}$, iterations of individual neurons can be vastly different, depends on their representation scheme $\mathcal{R}$ for individual sub-description of subcomponent - for example, if there exist the component called the \textit{loss function}, then the macroscopic behaviour will only be configured up to certain accuracy, the rest is left for the detail microscopic setting. Hence, we can partition $\mathcal{N}$ into $$\mathcal{N}\supset \mathcal{N}_{\mathcal{H}}(\dots)\times \mathcal{N}_{\mathcal{P}}(\dots), \quad \mathcal{N}_{M}\supset \mathcal{N}_{\mathcal{H},M}(n_{i},n_{o},M)\times \mathcal{N}_{\mathcal{P},M}(n_{i},n_{o},M)$$
of which here we use the word \textit{hyperparameter} for $\mathcal{H}$ a bit different from its usage of literature, and $\mathcal{P}$ for microscopic \textit{parameters}, under the proposition that the properties can be configured using parameters as abstraction. Similar notion goes for $\mathcal{C}$, though it is less apparent. 

Therefore, the minimally defined neuron, or what we can take as the \textbf{fundamental model of neural processing unit} works well, and provides us with the first iteration of the constructing brick for our larger, wider implementation. Though our aim is to advance this structure, first, we must take a look to the formalization and propositional-nation of the rudimentary concept of the neuron, even of the fundamental, minimal neuron. And this will happen at the end of this chapter. Until then, examining the maximal usage of the current structure will be more than enough as it aligns with historical literature's neuron and neural network formalism. Following sections will see how far we can go with this. 
\subsection{Why we call it minimal}

It seems a bit weird why we call our construction on its own a \textit{minimal construction}. To see this, though, we have to go back to the definition of a neuron, and the historical viewpoint on the constructing principle. The need for structuralization and complexity measure is impending on the theory itself, for example, the fact that there exists no current singular class of measure for any specific structures of feedforward networks, as seen of \cite{hu_model_2021}. Often, the minimal construction in typical layered network is the neuron unit only, of which is simply understood as a function with changing activation function. However, this is a limited approach, as it only satisfy mathematical elegance or simplicity of form, which it does well on standard systems, without targeting specific automaton-reliant and actual ergodic properties. For example, take the polynomial neuron. By certain structural minimization as we defined of the fundamental theorem

\section{Formalization}

The structure of our theory on the neural formalism is influenced by the object-abstracted treatment of mathematically embedded structures, and the unit-wise principle of particular neuron. Before we meet ourselves into the notion of epistemic circularity\footnote{Also called bootstrapping, where to understand or define certain notion, requires the knowledge of the object wishes to be defined itself - for example, defining the size of a finite natural number set, using the elements of the number set itself.} problem, we might as well clarify a few prerequisites for such structure to exhibit. 

First, we indict on the fundamental encoding environment that any object can take. The main point of any structure here is that there exists fundamentally the encoding space of two types. First is the object's cardinality space, denoted $\Gamma=(\mathbb{N},F)$ for any given categorization $F$. Second is the encoding primitive of the field $\mathbb{R}$ for generality - in general any field is alright, and they define the analytic structure of the system itself. Any extension, for example, the $\mathbb{R}$-algebra of complex number $\mathbb{C}$ is then the primitive field's extension. \footnote{This corresponds to a variety of idea. For example, see \cite{Cartuyvels_2021} for the idea of discrete and continuous representations and processing, and \cite{M_ller_2022} for the same idea, but used on computer graphic process where it is constructed as discrete indexing structure (hash tables) with continuous neural fields. Cardinality space takes inspiration from combinatorial species, sheaf theory of global-local set, and Grothendieck's approach to algebraic topology. Extension is the idea from field extension itself, for structures that can be extended. The general idea of this section relies on interpretation of category theory.}. Any type of data or system can then be decomposed to such, with additional structure on top of such primitive. Such is then called the \textit{primitive framework}. 

\begin{definition}[Primitive framework]
    Let us define the primitive framework $\mathcal{P}_{0}$ of the dual $(\Gamma, \mathbb{R})$ where $\Gamma=(\mathbb{N},F)$ is the cardinality encoding space, and $\mathbb{R}$ is the base field primitive of the analytical encoding. An object $X\in \mathcal{P}_{0}$ admits a dual representation, 
    \begin{equation}
        X \mapsto (\gamma(X),\rho(X))
    \end{equation}
    for $\gamma: \mathrm{Obj}\to \Gamma$ of cardinality encoding, and $\rho:\mathrm{Obj}\to\mathcal{E}(\mathbb{R})$ for the field extension of $\mathbb{R}$ category, or the category of all $\mathbb{R}$-algebras. For $\mathcal{E}(\mathbb{R})$ without extension, then $\mathcal{E}(\mathbb{R})\cong \{\mathbb{R}\}$ of all $\mathbb{R}$-algebra. 
\end{definition}

For a structure that is supposed to be unit-wise constructed like neural network and the like, one of the main principle is the principle of abstraction. With this, come the idea of layer. Specifically, a \textit{layer} separates abstraction in terms of subspace. Let us take an example of such kind for clarification. Let us define the primitive framework $\mathcal{P}_{0}$ as now the layer $L_{0}$ of this layering scheme. Then, we define the layer $\mathcal{L}_{1}$ of all unit-wised neuron-like units taking over the representation scheme on $\mathcal{L}_{0}$. We then define the structure of the neuron class $\mathcal{N}_{1}$ upon such as followed. 

\begin{definition}[Base neuron class]
    We define the base neuron class $\mathcal{N}_{1}\equiv L_{1}$ as followed. For any $\mathcal{U}\in L_{1}$ for $\mathcal{U}$ as unit-wise construction over $L_{0}$, then every $\mathcal{U}$ satisfies the input signature $\mathcal{I}_{sig}:\mathcal{E}'(\mathbb{R})\to L_{0}$, output signature $\mathcal{O}_{sig}:L_{0}\to\mathcal{E}'(\mathbb{R})$, and $\mathcal{C}_{ext}$ as extensible construction over $L_{0}$, for $\mathcal{E}'(\mathbb{R})$ particular extension on analytical encoding of $L_{0}$. The \textbf{type} template of a neural unit $\mathcal{U}\in L_{0}\equiv \mathcal{N}_{1}$ is then defined as $\mathcal{U}=\langle \mathcal{I}_{sig},\mathcal{C}_{ext},\mathcal{O}_{sig}\rangle$ where $\mathcal{I}_{sig}$ and $\mathcal{O}_{sig}$ are invariant as type. 
\end{definition}

Those definitions directly link it to type theory, while such development is perhaps more complex. In general, this specifies the \textit{type} of a particular unit using the invariant analytical typing in the general environment linking to it. This is the \textit{outer typing} of a given model construct \index{outer typing}, of which defers the type of which interaction between the environment, or the \textit{global space state} happens and what is received of such. While the ambiance space can be forgiving, such cannot be said for the typing of given structure, since it must have the correct typing for identification. The extension is then is fairly simple, but the more important notion is the equivalence of the variant $C_{ext}$ extensible unit of the inner structure itself. Hence, the equivalence will be of the following type: 



Let's take an example on it. The first example that we can think of refers to the multiple input neurons previously introduced. How can it be expressed under this framework? 

